\centerline{\bf \Huge Сириусовский счёт в комплах}

\begin{flushright}
        \scriptsize{бээээээм... домашний}
\end{flushright}

\problem{8.2}
Дан прямоугольный треугольник $АВС$ с острым углом $С$. На гипотенузе $АВ$ отмечена середина $L$. На стороне $BC$ отмечена точка $N$ такая, что $CN:NB = 1:2.$ Докажите, что $\angle CLN = \angle BAN$

\begin{proof}
    Введём комплексную систему координат так, что $C$ --- начало координат (то есть 0), $CA$ --- действительная ось, $CB$ --- мнимая ось. Тогда вершине $A$ соответствует комплексное число $a$, которое чисто действительное (то есть без мнимой части), а вершине $B$ соответствует комплексное число $bi$, которое чисто мнимое (то есть без действительной части). Здесь $a, b \in \mathbb{R}$ --- действительные числа. 
    
    Найдём комплексные координаты точек:

    $$L = \dfrac{A+B}{2} = \dfrac{a+bi}{2}$$

    $$N = \dfrac{2\cdot C + B}{3} = \dfrac{0 + bi}{3} = \dfrac{bi}{3}$$

    Найдём вектора, образующие интересующие нас углы:

    $\overrightarrow{C L}=L-C=\dfrac{a+b i}{2}$

    $\overrightarrow{L N}=N-L=\dfrac{b i}{3}-\dfrac{a+b i}{2}=-\dfrac{a}{2}-\dfrac{b i}{6}$

    $\overrightarrow{B A}=A-B=a-b i$

    $\overrightarrow{A N}=N-A=\dfrac{b i}{3}-a$

Угол $\angle C L N$ - это угол между векторами $\overrightarrow{C L}$ и $\overrightarrow{L N}$ :

$$
\theta_1=\arg \left(\dfrac{\overrightarrow{L N}}{\overrightarrow{C L}}\right)=\arg \left(\dfrac{-\dfrac{a}{2}-\dfrac{b i}{6}}{\dfrac{a+b i}{2}}\right)=\arg \left(\dfrac{-3 a-b i}{3 a+3 b i}\right)
$$


\newpage

\hspace{3mm}

Угол $\angle B A N$ - это угол между векторами $\overrightarrow{B A}$ и $\overrightarrow{A N}$ :

$$
\theta_2=\arg \left(\dfrac{\overrightarrow{A N}}{\overrightarrow{B A}}\right)=\arg \left(\dfrac{\dfrac{b i}{3}-a}{a-b i}\right)=\arg \left(\dfrac{-3 a+b i}{3 a-3 b i}\right)
$$

Заметим, что числа $\dfrac{-3a - bi}{3a + 3bi}$ и $\dfrac{- 3a + b i}{3a - 3b i}$ являются комплексно сопряженными, поэтому их аргументы равны по модулю (они будут разными по знаку, но в данном случае это не имеет значения). Что и требовалось доказать.

\end{proof}